% Part: second-order-logic
% Chapter: sol-and-set-theory
% Section: power-of-continuum

\documentclass[../../../include/open-logic-section]{subfiles}

\begin{document}

\olfileid{sol}{set}{pow}

\olsection{قوة المتصل}

\begin{explain}
نستطيع في منطق الرتبة الثانية أن نكمم على المجموعات الجزئية من
الـ!!{domain}، لكن لا نستطيع أن نكمم على مجموعات المجموعات الجزئية من
الـ!!{domain}.
ولفعل ذلك مباشرة نحتاج إلى منطق \emph{الرتبة الثالثة}. فلو أردنا، مثلًا،
تقرير مبرهنة كانتور القائلة إنه لا توجد دالة~!!{injective} من مجموعة
قوى مجموعة ما إلى المجموعة نفسها، فقد نحاول صياغتها هكذا: «لكل
مجموعة~$X$ ولكل مجموعة~$P$، إذا كانت~$P$ مجموعة قوى~$X$، فليس
$\cardle{P}{X}$». والقول إن~$P$ مجموعة قوى~$X$ يقتضي أن نصوغ أن
الـ!!{element}s في~$P$ هي جميع المجموعات الجزئية من~$X$ ولا شيء
سواها، أي شيئًا من قبيل~$\lforall[Y][(P(Y)
  \liff Y \subseteq X)]$. وتكمن المشكلة في~$P(Y)$: فهذه ليست
!!a{formula} من منطق الرتبة الثانية، لأن وسائط متغيرات العلاقة الأحادية
الموضع مثل~$P$ لا يمكن أن تكون إلا حدودًا.

غير أننا نستطيع \emph{محاكاة} التكميم على مجموعات المجموعات إذا كان
المجال كبيرًا بما يكفي. والفكرة هي الاستفادة من أن العلاقة الثنائية
الموضع~$R$ تربط~!!{element}s من الـ!!{domain} بـ!!{element}s من
الـ!!{domain}. فإذا أعطيت علاقة~$R$ من هذا القبيل، أمكننا جمع جميع
الـ!!{element}s التي يرتبط بها~$x$ ما بالعلاقة~$R$:
$\Setabs{y \in \Domain{M}}{R(x,y)}$ هي المجموعة التي «يرمزها»~$x$.
وبالعكس، إذا كانت~$Z \subseteq \Pow{\Domain{M}}$ مجموعة ما من
المجموعات الجزئية من~$\Domain{M}$، وكان في~$\Domain{M}$ من
الـ!!{element}s عدد لا يقل عن عدد المجموعات في~$Z$، فثمة أيضًا
علاقة~$R \subseteq \Domain{M}^2$ يرمز بالعلاقة~$R$ عنصر~$x$ ما
كل~$Y \in Z$.
\end{explain}

\begin{defn}
إذا كانت~$R \subseteq \Domain{M}^2$، فإن~\emph{$x$ يرمز بالعلاقة~$R$}
المجموعة~$\Setabs{y \in \Domain{M}}{R(x, y)}$.
\end{defn}

إذا كان~!!a{element}~$x \in \Domain{M}$ يرمز بالعلاقة~$R$ مجموعة
$Z \subseteq \Domain{M}$، فإن مجموعة~$Y \subseteq \Domain{M}$ ترمز
مجموعة من المجموعات، هي المجموعات التي ترمزها الـ!!{element}s من~$Y$.
ومن ثم يمكن لـ$Y$ أن ترمز بالعلاقة~$R$ المجموعة~$\Pow{X}$. وتفعل ذلك
إذا وفقط إذا كان، لكل~$Z \subseteq X$، عنصر~$x \in Y$ يرمز
بالعلاقة~$R$ المجموعة~$Z$، وكان كل~$x \in Y$ يرمز بالعلاقة~$R$
مجموعةً~$Z \subseteq X$.

\begin{prop}
تعبّر الـ!!{formula}
  \[
  \fn{Codes}(x, R, Z) \ident \lforall[y][(Z(y) \liff
    R(x, y))]
  \]
عن أن~$s(x)$ يرمز بالعلاقة~$s(R)$ المجموعة~$s(Z)$. وتعبّر
الـ!!{formula}
\begin{multline*}
  \fn{Pow}(Y, R, X) \ident \\
  \lforall[Z][(Z \subseteq X \lif \lexists[x][(Y(x) \land
    \fn{Codes}(x, R, Z))])] \land {} \\ \lforall[x][(Y(x) \lif
  \lforall[Z][(\fn{Codes}(x, R, Z) \lif Z \subseteq X)])]
\end{multline*}
عن أن~$s(Y)$ ترمز بالعلاقة~$s(R)$ مجموعة القوى لـ$s(X)$؛ أي إن عناصر
$s(Y)$ ترمز بالعلاقة~$s(R)$ جميع المجموعات الجزئية من~$s(X)$ ولا شيء
سواها.
\end{prop}

\begin{explain}
بهذه الحيلة نستطيع التعبير عن عبارات تخص مجموعة القوى بالتكميم على
رموز المجموعات الجزئية بدلًا من المجموعات الجزئية نفسها. فيمكن، مثلًا،
التعبير الآن عن مبرهنة كانتور بالقول إنه لا توجد دالة~!!{injective} من
مجال أي علاقة ترمز مجموعة القوى لـ$X$ إلى~$X$ نفسها.
\end{explain}

\begin{prop}
الجملة
\begin{multline*}
  \lforall[X][\lforall[Y][\lforall[R][(\fn{Pow}(Y, R, X) \lif \\
      \lnot \lexists[u][(\lforall[x][\lforall[y][(\eq[u(x)][u(y)] \lif
            \eq[x][y])]] \land {}\\
        \lforall[x][(Y(x) \lif X(u(x)))])])]]]
\end{multline*}
صادقة منطقيًا.
\end{prop}

\begin{explain}
مجموعة قوى مجموعة~!!a{denumerable} تكون~!!{nonenumerable}، ومن ثم
تكون أصليتها أكبر من أصلية أي مجموعة~!!{denumerable} (وهي
$\aleph_0$). ويسمى حجم~$\Pow{\Nat}$ «قوة المتصل»، لأنه يساوي حجم نقاط
خط الأعداد الحقيقية~$\Real$. وإذا كان الـ!!{domain} كبيرًا بما يكفي
لترميز مجموعة قوى مجموعة~!!a{denumerable}، أمكننا التعبير عن كون
مجموعة ما من حجم المتصل بالقول إنها مكافئة عدديًا لأي مجموعة~$Y$ ترمز
مجموعة قوى مجموعة~$X$ من الحجم~$\aleph_0$. (وإذا لم يكن المجال كبيرًا
بما يكفي، أي إذا لم يحتو مجموعة جزئية مكافئة عدديًا لـ$\Real$، فلن
توجد أيضًا علاقة ترمز~$\Pow{X}$.)
\end{explain}

\begin{prop}
إذا كان~$\cardle{\Real}{\Domain{M}}$، فإن الـ!!{formula}
\begin{multline*}
\fn{Cont}(Y) \ident 
\lexists[X][\lexists[R][((\fn{Aleph}_0(X) \land
      \fn{Pow}(Y, R, X)) \land \\ 
      \lforall[x][\lforall[y][((Y(x) \land {} 
      Y(y) \land \lforall[z][R(x,z) \liff R(y,z)]) \lif \eq[x][y])]])]]
\end{multline*}
تعبّر عن أن~$\cardeq{s(Y)}{\Real}$.
\end{prop}

\begin{proof}
تعبّر~$\fn{Pow}(Y, R, X)$ عن أن~$s(Y)$ ترمز بالعلاقة~$s(R)$ مجموعة
القوى لـ$s(X)$، التي تقول~$\fn{Aleph}_0(X)$ إنها قابلة للتعداد. ومن
ثم تكون~$s(Y)$ كبيرة بقدر قوة المتصل على الأقل، وإن جاز أن تكون أكبر
(إذا كانت عناصر متعددة من~$s(Y)$ ترمز المجموعة الجزئية نفسها من~$X$).
ويستبعد الحد الأخير هذا الاحتمال، إذ يقتضي أن يكون الربط، بالعلاقة
$s(R)$، بين عناصر~$s(Y)$ والمجموعات الجزئية من~$s(X)$~!!{injective}.
\end{proof}

\begin{prop}
تكون~$\cardeq{\Domain{M}}{\Real}$ إذا وفقط إذا
\begin{multline*}
  \Sat{M}{\lexists[X][\lexists[Y][\lexists[R][(\fn{Aleph}_0(X) \land \fn{Pow}(Y, R, X) \land {}\\
      \lforall[x][\lforall[y][((Y(x) \land Y(y) \land
        \lforall[z][R(x,z) \liff R(y,z)]) \lif \eq[x][y])]] \land {}\\
      \lexists[u][(\lforall[x][\lforall[y][(\eq[u(x)][u(y)] \lif \eq[x][y])]] \land {}
        \\\lforall[x][Y(u(x))] \land {}
        \\\lforall[y][(Y(y) \lif \lexists[x][\eq[y][u(x)]])])])]]]}. 
\end{multline*}
\end{prop}

\begin{explain}
فرضية المتصل هي القول إن حجم المتصل هو الأصلية الأولى
!!{nonenumerable}؛ أي إن~$\Pow{\Nat}$ من الحجم~$\aleph_1$.
\end{explain}

\begin{prop}
تصدق فرضية المتصل إذا وفقط إذا كانت
\[\fn{CH} \ident
\lforall[X][(\fn{Aleph}_1(X) \liff \fn{Cont}(X))]\]
صادقة منطقيًا.
\end{prop}

لاحظ أنه لا يصدق أن~$\lnot \fn{CH}$ صادقة منطقيًا إذا وفقط إذا كانت
فرضية المتصل كاذبة. ففي بنية~!!a{enumerable} لا توجد مجموعات جزئية من
الحجم~$\aleph_1$، ولا مجموعات جزئية من حجم المتصل؛ ومن ثم تكون
$\fn{CH}$ صادقة دائمًا في بنية~!!a{enumerable}. غير أننا نستطيع إعطاء
جملة مختلفة تصدق منطقيًا إذا وفقط إذا كانت فرضية المتصل كاذبة:

\begin{prop}
تكون فرضية المتصل كاذبة إذا وفقط إذا كانت
\[\fn{NCH} \ident
\lforall[X][(\fn{Cont}(X) \lif \lexists[Y][(Y \subseteq X \land \lnot
    \fn{Count}(Y) \land \lnot \cardeq{X}{Y})])]\]
صادقة منطقيًا.
\end{prop}

\end{document}
